\documentclass[addpoints,answers,12pt]{exam}
\usepackage{amsmath,amssymb,mathtools}
\usepackage{bm}
\usepackage{graphicx}
\usepackage{enumitem}
\usepackage{tikz}
\usepackage{geometry}
\geometry{margin=1in}

\pagestyle{headandfoot}
\firstpageheader{ECON 5001 (Healy)}{Midterm Exam}{October 28, 2026}
\runningheader{ECON 5001 (Healy)}{Midterm Exam Page \thepage\ of \numpages}{10/28/2026}
\firstpagefooter{}{}{}
\runningfooter{}{}{}

\begin{document}

\CorrectChoiceEmphasis{\bfseries \color{red}}

\vspace*{0.1in}

\makebox[\textwidth]{Name:\enspace\hrulefill}

\vspace{0.2in}

\textbf{Instructions:} Please write your name! This exam has \numquestions\ questions given over \numpages\ pages, including blank pages for scrap work at the end. Please verify you have all pages before you begin. The total number of points on this exam is \numpoints. For multiple choice questions, please write your answer clearly in the \textbf{left margin} of each question. If you need extra scrap paper, please see the proctor.

\vspace{0.2in}

\section*{Multiple Choice Questions:}

\begin{questions}

%%%%% Multiple choice: 26 items %%%%%

\question[1] Let $f(x) = -2x^2 + 4x + 1$. When searching for a maximum or minimum, which is true?
\begin{choices}
    \choice $f'(x)=0$ at $x=2$, which is a local maximum since $f''(2)>0$.
    \choice $f'(x)=0$ at $x=1$, which is a local minimum since $f''(1)>0$.
    \CorrectChoice $f'(x)=0$ at $x=1$, which is a local maximum since $f''(1)<0$.
    \choice This function has no maximum or minimum.
\end{choices}

\question[1] Use the Lagrangian technique to maximize $u(x,y)=\ln(x) + 2\ln(y)$ subject to $p_x x + p_y y = I$. (Recall that the derivative of $\ln(x)$ is $1/x$ and the derivative of $2\ln(y)$ is $2/y$.) The resulting demand functions are:
\begin{choices}
    \choice $x = \dfrac{I}{p_x + p_y},\; y=\dfrac{I}{p_x + p_y}$.
    \choice $x = \dfrac{2I}{3p_x},\; y=\dfrac{I}{3p_y}$.
    \CorrectChoice $x = \dfrac{I}{3p_x},\; y=\dfrac{2I}{3p_y}$.
    \choice $x = \dfrac{I}{p_x},\; y=0$.
\end{choices}

\question[1] Consider $f(x,y)=x^2+y-1$. Let $y(x)$ be the implicit function that solves $x^2 + y(x) -1 =0$ for any $x$. What are $y(x)$, $y'(x)$, and $\dfrac{-\partial f/\partial x}{\partial f/\partial y}$?
\begin{choices}
    \choice $y(x)=x^2+y-1$, $y'(x)=2x$, and $(-\partial f/\partial x)/(\partial f/\partial y)=-2x$.
    \CorrectChoice $y(x)=1-x^2$, $y'(x)=-2x$, and $(-\partial f/\partial x)/(\partial f/\partial y)=-2x$.
    \choice $y(x)=x^2-1$, $y'(x)=2x$, and $(-\partial f/\partial x)/(\partial f/\partial y)=2x$.
    \choice The implicit function cannot be solved.
\end{choices}

% \question[1] Let a consumer maximize $u(x,y)=x^{1/2}y^{1/2}$ subject to $x+2y=100$. In the Lagrangian, the first two equations $(\partial \mathcal{L}/\partial x$ and $\partial\mathcal{L}/\partial y$) tell us that:
% \begin{choices}
%     \choice $y = x$.
%     \choice $y = 2x$.
%     \CorrectChoice $x = 2y$.
%     \choice Corner solution $x=100,y=0$.
% \end{choices}

\question[1] A farmer with 100 ft of fence wants to fence off a rectangular field with one side against an existing straight barn wall (no fencing needed on that side). If $x$ is the length parallel to the wall and $y$ the length perpendicular to the wall, then the area is $f(x,y)=xy$ and the constraint is $x + 2y = 100$. The area-maximizing dimensions satisfy
\begin{choices}
    \choice $x=25, y=25$.
    \choice $x=50, y=25$.
    \CorrectChoice $x=50, y=25$.
    \choice $x=33.\overline{3}, y=33.\overline{3}$.
\end{choices}

\question[1] Recall that a production function $f(k,l)$ is homogenous of degree $r$ if\\
$f(\alpha k,\alpha l)=\alpha^r f(k,l)$. When $r=1$ we say that production has
\begin{choices}
    \choice decreasing returns to scale (DRS).
    \CorrectChoice constant returns to scale (CRS).
    \choice increasing returns to scale (IRS).
    \choice None of the above; homogeneity degree is not related to returns to scale.
\end{choices}

\question[1] Which axiom says that if $a\succsim b$ and $b\succsim c$ then $a\succsim c$?
\begin{choices}
    \choice Completeness.
    \CorrectChoice Transitivity.
    \choice Continuity
    \choice Quasiconcavity.
\end{choices}

\question[1] If both $x$ and $y$ are ``goods'' (meaning the consumer likes more of each) then the indifference curves must be
\begin{choices}
    \choice inverted.
    \choice upward sloping.
    \CorrectChoice downward sloping.
    \choice cannot be determined without further information.
\end{choices}

\question[1] For $u(x,y)=x+y$, indifference curves are
\begin{choices}
    \choice strictly convex.
    \choice strictly concave.
    \CorrectChoice linear.
    \choice upward sloping.
\end{choices}

\newpage
\question[1] If goods are perfect complements with utility $u(x,y)=\min\{x,y\}$, the marginal rate of substitution (MRS) is
\begin{choices}
    \CorrectChoice infinite above the kink point, zero right of the kink point, and undefined at the kink point.
    \choice undefined everywhere.
    \choice infinite everywhere.
    \choice zero everywhere.
\end{choices}

\question[1] A consumer with Cobb-Douglas utility $u(x,y)=x^{\alpha}y^{1-\alpha}$ facing prices $p_x,p_y$ and income $I$ chooses $x$ such that
\begin{choices}
    \CorrectChoice $x = \alpha I/p_x$.
    \choice $x = \alpha I/p_y$.
    \choice $x = (1-\alpha)I/p_x$.
    \choice $x = I/(p_x + p_y)$.
\end{choices}

\question[1] With two goods $x$ and $y$, the endpoints of the consumer's budget line are given by
\begin{choices}
    \choice $x=I-p_x$ and $y=I-p_y$.
    \choice $x=I-p_y$ and $y=I-p_x$.
    \choice $x=I/p_y$ and $y=I/p_x$.
    \CorrectChoice $x=I/p_x$ and $y=I/p_y$.
\end{choices}

\question[1] Consider the utility maximization problem $\max u(x,y)$ s.t. $p_x x+p_y y=I$. Two cases where the Lagrangian method may fail are when
\begin{choices}
    \choice $u(x,y)=x^\alpha y^{(1-\alpha)}$ and $u(x,y)=\min\{x,2y\}$.
    \CorrectChoice $u(x,y)=x+2y$ and $u(x,y)=\min\{x,2y\}$.
    \choice $u(x,y)=x^\alpha y^{(1-\alpha)}$ and $u(x,y)=x-y$ (because $y$ is a ``bad'').
    \choice $u(x,y)=x^\alpha (-y)^{(1-\alpha)}$ (because $y$ is a ``bad'') and $u(x,y)=2x+y$.
\end{choices}

\question[1] Suppose for good $x$ demand is $x^*(p_x,p_y,I)=\dfrac{I}{p_x}$. Income rises from $I$ to $2I$. Demand for $x$
\begin{choices}
    \choice grows by a factor of $2/p_x$.
    \CorrectChoice doubles.
    \choice halves.
    \choice remains unchanged.
\end{choices}

\newpage
\question[1] Suppose there's only one good, $x$, and utility is given by $u(x)=x^2$. The consumer faces budget constraint $p_x x = I$. If the price $p_x$ increases, the consumer's demand $x^*(p_x,I)$
\begin{choices}
    \choice will stay the same since their total expenditure on $x$ is fixed.
    \choice will increase since it's a Giffen good.
    \CorrectChoice will decrease since their total expenditure on $x$ is fixed.
    \choice is undefined because the Lagrangian gives $x=\lambda p_x/2$ but we don't know $\lambda$.
\end{choices}

\question[1] Two goods are gross substitutes if
\begin{choices}
    \choice an increase in the price of good 1 decreases demand for good 2.
    \CorrectChoice an increase in price of good 1 increases demand for good 2.
    \choice the price of good 1 has no affect on the demand for good 2.
    \choice an increase in income increases the demand for both goods.
\end{choices}

\question[1] For Leontief (fixed proportions) production $f(k,l)=\min\{\alpha l,\beta k\}$, the isoquant is
\begin{choices}
    \choice smooth and convex.
    \choice linear with slope $-\alpha/\beta$.
    \choice upward sloping.
    \CorrectChoice a right-angle (kinked).
\end{choices}

\question[1] For Leontief production $f(k,l)=\min\{\alpha l,\beta k\}$, the cost to produce $q$ units (at the lowest cost) is
\begin{choices}
    \choice $C(q) = \left(v^\alpha w^{1-\alpha}\frac{1}{\alpha^\alpha(1-\alpha)^{1-\alpha}}\right)q$.
    \CorrectChoice $C(q) = \left(\frac{w}{\alpha}+\frac{v}{\beta}\right) q$.
    \choice $C(q)=\max\{\frac{w}{\alpha},\frac{v}{\beta}\}\cdot q$.
    \choice quadratic in $q$.
\end{choices}

\question[1] If a firm's cost function is linear, meaning $C(q)=FC + aq$ where $FC$ and $a$ are numbers, then the marginal cost function is
\begin{choices}
    \choice $MC(q)=FC$ (constant).
    \choice $MC(q)=aq + FC$ (linear).
    \choice $MC(q)=aq$ (linear).
    \CorrectChoice $MC(q)=a$ (constant).
\end{choices}

\newpage
\question[1] Recall that a monopolist who faces prices $p(q)=100-2q$ and constant marginal cost $MC=20$ chooses $q$ to satisfy $MR(q)=MC(q)$. The profit-maximizing output (if they produce) is therefore
\begin{choices}
    \choice $q=0$, since $MR(q)=MC(q)$ gives ``$-2=20$'', so revenue is always negative.
    \CorrectChoice $q=20$.
    \choice $q=40$.
    \choice $q=50$.
\end{choices}

\question[1] A competitive (price-taking) firm has total cost $C(q)=2q^3 -40 q^2 + 250q$, with no fixed costs. Thus, $VC(q)=2q^3-40q^2+250q$ and $FC=0$. Which of the following describes their short-run shutdown decision and profit levels?
\begin{choices}
    \choice If $p\ge 10$ then they operate at a profit. Otherwise, shut down.
    \CorrectChoice If $p\ge 50$ then they operate at a profit. Otherwise, shut down.
    \choice If $0\le p\le 10$ then they operate at a loss, and if $p\ge 10$ then they operate at a profit.
    \choice If $p< 10$ they shut down. If $10\le p\le 50$ they operate at a loss. If $p\ge 50$ they operate at a profit.
\end{choices}

\question[1] Consider a partial equilibrium setting with $100$ consumers. If individual demand for good $x$ is given by $x_i^*(p_x,p_y,I)=20 - 2p_x$ for each consumer $i$, then the total market demand is
\begin{choices}
    \choice $Q_D(p_x,p_y,I)=20-2p_x$.
    \choice $Q_D(p_x,p_y,I)=2000 - 2 p_x$.
    \CorrectChoice $Q_D(p_x,p_y,I)=2000-200p_x$.
    \choice $Q_D(p_x,p_y,I)=20-2(100p_x)$.
\end{choices}

\question[1] If entry and exit doesn't affect costs then long-run supply is
\begin{choices}
    \choice upward sloping.
    \CorrectChoice horizontal at $p=\min AC$.
    \choice horizontal at $p=\min AVC$.
    \choice downward sloping.
\end{choices}

\question[1] The Income Identity says that
\begin{choices}
    \choice all firms act as if they're just maximizing revenue, since total costs in the world are fixed.
    \choice all consumers have the same income $I$.
    \choice if we double everyone's income and double all prices, demand for all goods will double.
    \CorrectChoice the total income of all people in the world equals the total revenue of all firms in the world.
\end{choices}

\question[1] In the Edgeworth box for production with two goods $x$ and $y$ produced from inputs $l$ and $k$, an efficient input combination requires
\begin{choices}
    \choice equal utilities of all consumers.
    \choice equal marginal rates of substitution of the consumers ($MRS_x=MRS_y$).
    \choice equal levels of $x$ and $y$.
    \CorrectChoice equal slopes of the isoquants ($MRTS_x=MRTS_y$).
\end{choices}

\question[1] The First Welfare Theorem says that the competitive equilibrium allocation will equal the social optimum, meaning that a benevolent social planner would pick the same allocation as the free market. Which assumption was \textit{not} used to derive this result?
\begin{choices}
    \CorrectChoice Prices are set by one large firm acting as a monopolist.
    \choice Complete, transitive, and continuous preferences.
    \choice A large number of firms and consumers who all take prices as given.
    \choice All firms and consumers have complete information about the goods being sold.
\end{choices}

%%%%% End MC questions %%%%%

\newpage

\section*{Short Answer/Calculation Questions:}

\question[10]
A consumer has preferences represented by $u(x,y)=x^{1/3} y^{2/3}$.
\begin{parts}
    \part[4] What are their demand functions $x(p_x,p_y,I)$ and $y(p_x,p_y,I)$?
    \part[3] Are these goods gross substitutes, gross complements, or neither?
    \part[3] Now suppose the price of $x$ is $p_x=2$, the price of $y$ is $p_y=8$, and income $I=600$. Compute the actual demand levels for $x$ and $y$ at these prices and income.
\end{parts}

\vspace{1em}
\begin{solution}[6cm]
For Cobb-Douglas with exponents $\alpha=1/3$ on $x$ and $1-\alpha=2/3$ on $y$, expenditure shares are $\alpha$ and $1-\alpha$. Thus $x=\dfrac{\alpha I}{p_x}=\dfrac{(1/3)I}{p_x},\; y=\dfrac{(2/3)I}{p_y}$. Plug numbers: $x=\dfrac{200}{2}=100$, $y=\dfrac{400}{8}=50$. Demand for $x$ doesn't depend on $p_y$ and demand for $y$ doesn't depend on $p_x$, so they are neither complements or substitutes.
\end{solution}

\newpage
\question[10]
Consider a price-taking firm with short-run cost function $C(q)=q^2 + 6q + 4$. (Notice the fixed cost of $4$.) Output price is $p$.
\begin{parts}
    \part[4] Derive the firm's short-run supply function (including shutdown condition).
    \part[3] Suppose $p=10$. Compute the profit-maximizing $q$ and the firm's profit.
    \part[3] For what prices will the firm shut-down?
\end{parts}

\vspace{1em}
\begin{solution}[6cm]
(a) $MC=6+2q$. For competitive firm set $p=MC$ so $q=(p-6)/2$ (if $p\ge 6$). Shutdown condition: produce if $p\ge \min AVC$. $AVC(q)=6 + q$ so $\min AVC$ is at $q=0$, which gives $\min AVC = 6$. So, produce if $p\ge 6$ and shutdown if $p<6$. Putting it together, the the supply function is $q=(p-6)/2$ if $p\geq 6$ and $q=0$ if $p<6$. Or, if you want to express price as a function of quantity (as it is in the usual graph), $p=2q+6$ for all $q\geq 0$. There's no actual discontinuity because the shut-down point happens exactly when $q=0$. (b) For $p=10$, $q=(10-6)/2=2$, profit $=pq - C(q)=10\cdot2 - (4+6\cdot2 +2^2)=20 - (4+12+4)=0$. (c) See above.
\end{solution}

\newpage
\question[10] Consider a General Equilibrium (GE) setting with 2 goods, $x$ and $y$.
\begin{parts}
    \part[2] Draw a Production Possibilities Frontier (PPF) that is curved downwards (concave, due to diminishing returns).
    \part[2] Next, draw the line representing both firms' total revenue and the representative consumer's budget at current prices $p_x$ and $p_y$. Identify the supply point.
    \part[2] Next, draw an indifference curve for the representative consumer and identify the demand point. Make it such that there is excess supply for $x$ and excess demand for $y$.
    \part[2] In this situation, how will prices change? Which way will the revenue/budget line rotate?
    \part[2] Finally, draw a new revenue/budget line after prices equilibrate so that supply=demand. Add another indifference curve for the consumer to show that their demand point matches the supply point.
\end{parts}

\vspace{1em}
\begin{solution}[6cm]
\\
%\includegraphics[width=6in]{graphics/GE_Equilibrium.png}\\
(Graph omitted.) The PPF is the boundary of the red area. The revenue/budget line is in blue. Supply is where that is tangent to the PPF. Demand is up to the left, giving excess supply for $x$ and excess demand for $y$. This means $p_x\downarrow$ and $p_y\uparrow$, so the revenue/budget line will rotate counterclockwise until we reach equilibrium, given by the purple revenue/budget line. For that line the supply point (tangecy with PPF) equals the demand point (tangency with the indifference curve).

\end{solution}

\newpage
\centering Page Left Blank for Scrap Work

\newpage
\centering Page Left Blank for Scrap Work

\newpage
\centering Page Left Blank for Scrap Work



\end{questions}

\end{document}
